\documentclass[10pt,journal,comsoc]{IEEEtran}
\usepackage{amsmath,amssymb,color}
\usepackage{amsfonts}
\usepackage{graphicx}
\usepackage[caption=false,font=footnotesize]{subfig}
\usepackage{cite}
\usepackage{url}
\usepackage{booktabs}
\usepackage{multirow}
\usepackage{makecell}

\begin{document}
% ============================
\section{Proposed Self-Interference-Assisted Monostatic Full-Duplex ISAC System}
\label{sec:system_model}

Fig.~\ref{fig:concept} shows the proposed monostatic full-duplex THz
ISAC system. The free-running photonic transmitter and the sensing
receiver share a single aperture through a circular-polarization
duplexer. The transmitted ISAC signal reflected from the communication
receiver, which serves as the sensing target, returns through the same
duplexer to the sensing receiver.

The difficulty of monostatic operation is that a portion of the strong
transmit signal leaks through the duplexer into the sensing receiver as
self-interference (SI), which is conventionally suppressed. The proposed
system instead uses it as a reference for direct square-law detection,
so that both links are realized with free-running lasers and envelope
detectors alone, without a local oscillator. Direct detection is
inherently phase-insensitive, but with subcarrier intensity modulation
its beat against the reference recovers the phase needed for both
communication and ranging.

Reusing the SI in this way serves two roles. As a strong,
range-independent reference it homodynes the weak echo, relaxing the
sensing output from an $R^{-8}$ to an $R^{-4}$ dependence; and as a
zero-delay reference plane it normalizes the estimated channel response,
removing the common gain and phase of the free-running source and
yielding an absolute, drift-immune range.

\begin{figure}[!t]
\centering
\includegraphics[width=\columnwidth]{fig/fig2.png}
\caption{Proposed self-interference-assisted monostatic full-duplex
THz ISAC system. LD: laser diode; MZM: Mach--Zehnder modulator; UTC-PD: uni-travelling-carrier
photodiode; LNA: low-noise amplifier; SLD: square-law detector; DSP:
digital signal processing; RHCP/LHCP: right-/left-hand circular
polarization.}
\label{fig:concept}
\end{figure}

\subsection{Transmit Signal Model}
\label{sec:tx_signal}

The transmitted THz waveform, normalized so that the carrier has unit
amplitude, is
\begin{equation}
x(t)=\big[1+m\,s(t)\big]\,
e^{j\left(2\pi (f_c+\delta\!f)t+\Delta\theta(t)\right)},
\label{eq:xt}
\end{equation}
where $m$ is the modulation index, $s(t)$ the real, unit-power
subcarrier signal carrying the symbol stream at the IF $f_{\mathrm{IF}}$,
and $f_c$ the beat frequency of the two lasers; the
transmitted power $P_t$ is carried by the link coefficients defined
below. Free-running operation introduces two impairments: the fast
common phase noise $\Delta\theta(t)$, and the slow carrier drift
$\delta\!f$, constant within one capture but unknown and varying across
captures. The carrier-to-sideband power ratio (CSPR) is $1/m^2$, whose choice
trades the desired beat term against nonlinear distortion in the
modulator and the photomixer. We adopt a conservative operating point,
$\mathrm{CSPR}=13$~dB and $m=0.22$, for the range analysis that
follows.

\subsubsection{Waveform}
The carrier drift precludes time-domain matched filtering, whose
correlation peak fades as $\delta\!f$ wanders; range must instead be
estimated from the channel frequency response (CFR), which is
drift-immune as shown below. The waveform is therefore drawn from the
orthogonal frequency-division multiplexing (OFDM) family, which pairs
naturally with per-subcarrier CFR estimation and frequency-domain
equalization. At THz frequencies the saturated output power of
amplifiers falls rapidly with carrier frequency, making a low
peak-to-average power ratio (PAPR) essential, and single-carrier
DFT-spread OFDM (DFT-s-OFDM) is accordingly regarded as a stronger THz
candidate than cyclic-prefix OFDM \cite{Han2023}. Its lower PAPR raises the
average transmit power attainable under the same peak-power limit.

\subsubsection{Power Allocation}
The waveform superimposes a deterministic sensing pilot and the
communication payload \cite{lyu_exploring_2025},
\begin{equation}
s(t)=\sqrt{\rho}\,p(t)+\sqrt{1-\rho}\,d(t),\qquad
\mathbb E\big[|s(t)|^2\big]=1,
\label{eq:st}
\end{equation}
where $p(t)$ is a unit-power constant-modulus pilot of $N_p$ symbols,
$d(t)$ the unit-power data symbols, and $\rho\in[0,1]$ the
power-allocation ratio. Both functions occupy the same time-frequency
resources, and $\rho$ trades the two links' maximum ranges.

\subsection{Sensing Branch}
\label{sec:sens_branch}

At the sensing port of the duplexer, which feeds the LNA input, the transmit
leakage superimposes on the target echo of round-trip delay
$\tau=2R/c$,
\begin{equation}
r(t)=\sqrt{P_{\mathrm{SI}}}\;x(t)+\sqrt{P_{\mathrm{ec}}(R)}\;x(t-\tau)
+n(t),
\label{eq:rt}
\end{equation}
where $n(t)$ is the input-referred receiver noise of power $N$.
The leakage power $P_{\mathrm{SI}}=P_t\,10^{-\mathrm{ISO}/10}$ is fixed
by the duplexer isolation and is independent of range, whereas the echo
power follows the radar equation \cite{sturm_waveform_2011},
\begin{equation}
P_{\mathrm{ec}}(\sigma_{\mathrm{eff}},R)
=\frac{P_tG_tG_r\lambda^2\sigma_{\mathrm{eff}}}{(4\pi)^3R^4},
\label{eq:pec}
\end{equation}
where $\sigma_{\mathrm{eff}}$ is the effective radar cross-section in
the observation direction, which absorbs the residual antenna mismatch,
the aspect-angle dependence, and the structural scattering of the
target. The power-allocation ratio $\rho$ does not enter here: it
divides the modulation sidebands of~\eqref{eq:st}, not the carrier that
dominates $P_t$, and therefore appears later in the SINR rather than in
the link budget. The duplexer isolation is chosen so that
$P_{\mathrm{SI}}<\mathrm{IP}_{1\mathrm{dB}}$ at the LNA input, which
sets the operating region of the front end.

\subsubsection{Square-Law Detection}
The amplified signal is rectified by a square-law detector (SLD),
implemented as a zero-bias Schottky diode, with output proportional to
its squared magnitude. Squaring the transmit waveform
of~\eqref{eq:xt} gives $|x|^2=1+2m\,s+m^2s^2$: DC blocking removes the
constant, the order-$m$ term is the wanted IF signal, and the
order-$m^2$ term is signal--signal beat interference (SSBI). The SSBI
spreads from DC over about twice the signal bandwidth, so only a
fraction $\kappa\le1$ of its power falls in the IF passband and
survives the subsequent processing. We carry this fraction explicitly
and write the in-band SSBI power as $\kappa m^4P^2$ for a path of power
$P$; the value in Table~\ref{tab:simulation_parameters} is a
conservative upper bound. Squaring~\eqref{eq:rt} and
grouping the products by role,
\begin{align}
|r(t)|^{2}=&\underbrace{2\sqrt{P_{\mathrm{SI}}P_{\mathrm{ec}}}\,
\mathrm{Re}\{x(t)x^{*}(t-\tau)\}}_{\text{SI--echo cross-beat}}
+\underbrace{P_{\mathrm{ec}}\,|x(t-\tau)|^{2}}_{\text{echo self-beat}}
\notag\\
&+\underbrace{P_{\mathrm{SI}}\,|x(t)|^{2}}_{\text{SI self-beat}}
+\mathcal N(t),
\label{eq:sq_expand}
\end{align}
where the noise term collects the beat of the leakage against the
receiver noise and the noise against itself,
\begin{equation}
\mathcal N(t)=2\sqrt{P_{\mathrm{SI}}}\,\mathrm{Re}\{x(t)n^*(t)\}
+|n(t)|^2.
\label{eq:noise}
\end{equation}
The first two terms of~\eqref{eq:sq_expand} carry the delay $\tau$ and
form the ranging signal. The SI--echo cross-beat arises only when the
leakage is present; without it, only the echo self-beat remains.

Every term in~\eqref{eq:sq_expand} is a product of a field with its own
conjugate, so the carrier phase and the phase noise $\Delta\theta(t)$
cancel and no local oscillator is required. Only the cross-beat retains
the delay-differential phase
$\Delta\phi_\tau=\Delta\theta(t)-\Delta\theta(t-\tau)$, whose effect is
small at short range: from
$\mathrm{var}(\Delta\phi_\tau)\simeq2\pi\Delta\nu\tau$, a linewidth
$\Delta\nu=15$~kHz and $\tau=13.3$~ns at $R=2$~m give
$\sigma_{\Delta\phi}\approx0.035$~rad.

Retaining the terms proportional to $m$, and at zero delay only the SI
self-beat, the sensing IF signal is
\begin{align}
y(t)=2m\big[&P_{\mathrm{SI}}\,s(t)\notag\\
&+\big(\sqrt{P_{\mathrm{SI}}P_{\mathrm{ec}}}\,e^{-j2\pi f_c\tau}
+P_{\mathrm{ec}}\big)s(t-\tau)\big]+n_{\mathrm{d}}(t),
\label{eq:yt}
\end{align}
the sum of an SI-borne copy at zero delay and an echo-borne copy at
delay $\tau$. The echo-borne amplitude comprises the cross-beat and the
echo self-beat, whose powers become the
$P_{\mathrm{SI}}P_{\mathrm{ec}}$ and $P_{\mathrm{ec}}^2$ terms of the
sensing SINR.

The detector-output noise $n_{\mathrm{d}}(t)$ has two parts: a floor
$N_{\mathrm{d},0}$ independent of the leakage, collecting the
noise--noise beat of~\eqref{eq:noise} together with the noise added
after detection by the diode and the IF amplifier, and the
leakage--noise beat, giving
$N_{\mathrm{d}}=N_{\mathrm{d},0}+2P_{\mathrm{SI}}N$. The leakage
therefore strengthens the sensing signal and raises the noise at the
same time. All powers are referred to the LNA input over the same noise
bandwidth, so this expression and the SINRs below are written in
equivalent input-referred power products rather than actual detector
output powers.

\subsubsection{SI-Referenced Range Estimation}
Fourier transforming~\eqref{eq:yt} and dividing by the known transmit
spectrum $S(f)$ yields the channel frequency response (CFR)
\begin{equation}
H(f)=A(f)e^{j\varphi}\Big[\sqrt{P_{\mathrm{SI}}}
+\sqrt{P_{\mathrm{ec}}}\;e^{-j2\pi(f+f_c)\tau}\Big],
\label{eq:cfr}
\end{equation}
where $A(f)e^{j\varphi}$ is common to both paths and collects the IF
channel response together with the phase left by the free-running
source, namely the per-capture offset due to the drift $\delta\!f$ and
the residual $\Delta\phi_\tau$.

The SI contribution is frequency-flat, whereas the echo oscillates with
period $1/\tau$, so averaging $H(f)$ across the band isolates the SI
path as $H_{\mathrm{SI}}(f)\approx
A(f)e^{j\varphi}\sqrt{P_{\mathrm{SI}}}$. Dividing by it gives the
normalized CFR
\begin{equation}
\tilde H(f)=\frac{H(f)}{H_{\mathrm{SI}}(f)}-1
=\sqrt{\frac{P_{\mathrm{ec}}}{P_{\mathrm{SI}}}}\;
e^{-j2\pi(f+f_c)\tau},
\label{eq:sinorm}
\end{equation}
which is the second role of the SI. For a single line-of-sight target
$\tilde H(f)$ is a pure linear phase, and the delay and range follow
from its slope,
\begin{equation}
\hat\tau=-\frac{1}{2\pi}\frac{\mathrm{d}\angle\tilde H(f)}{\mathrm{d}f},
\qquad \hat R=\frac{c}{2}\hat\tau.
\label{eq:slope}
\end{equation}

The normalization is what makes~\eqref{eq:slope} usable. Neither part
of $A(f)e^{j\varphi}$ is known in advance, and each corrupts the
estimate in its own way: the group delay of $\angle A(f)$ adds directly
to the measured slope and biases the range, while the source phase
changes from capture to capture. Dividing by $H_{\mathrm{SI}}$ removes
both, since they are common to the two paths, leaving a slope that is
the target delay alone, measured against the zero-delay SI path. The
drift survives only as a constant phase and does not affect the slope.
The range is therefore absolute and repeatable across captures without
external calibration.

For multiple targets the full range profile follows from an inverse
discrete Fourier transform of $\tilde H(f)$, and a target-free
reference capture may be subtracted to suppress static clutter
\cite{sturm_waveform_2011}.

\subsubsection{Sensing SINR}
Collecting from~\eqref{eq:sq_expand}--\eqref{eq:noise} the terms that
survive DC blocking and IF filtering, and applying the coherent
processing gain $G_p=N_p$ of the pilot correlation, the sensing
signal-to-interference-plus-noise ratio (SINR) is
\begin{equation}
\gamma_{\mathrm{sens}}=\rho\,G_p\,
\frac{2m^2\big[P_{\mathrm{SI}}P_{\mathrm{ec}}+P_{\mathrm{ec}}^2\big]}
{\underbrace{N_{\mathrm{d},0}}_{\text{detector floor}}
+\underbrace{2P_{\mathrm{SI}}N}_{\text{SI--noise}}
+\underbrace{\kappa m^4P_{\mathrm{SI}}^2}_{\text{SSBI}}},
\label{eq:sinr_sens}
\end{equation}
where $\rho$ enters as the pilot fraction of the modulation. The
$P_{\mathrm{SI}}P_{\mathrm{ec}}$ and $P_{\mathrm{ec}}^2$ terms of the
numerator are the cross-beat and the echo self-beat; the denominator
holds the detector-output floor, the SI--noise beat,
and the leakage SSBI.

The leakage strengthens the ranging signal through the cross-beat
$P_{\mathrm{SI}}P_{\mathrm{ec}}$, which scales as $R^{-4}$. The
SI--noise beat grows in proportion to the leakage, as the ranging term
does, and therefore saturates the SINR rather than reducing it, whereas
the SSBI grows as $P_{\mathrm{SI}}^2$ and eventually reverses the gain.

Setting $P_{\mathrm{SI}}=0$ leaves the echo detected through its own
self-beat alone against the fixed floor,
\begin{equation}
\gamma_{\mathrm{sens}}^{\mathrm{w/o\,SI}}
=\rho G_p\,\frac{2m^2P_{\mathrm{ec}}^2}{N_{\mathrm{d},0}}
\;\propto\;R^{-8},
\label{eq:sinr_nosi}
\end{equation}
where the echo SSBI is omitted because $P_{\mathrm{ec}}$ is small. The
maximum sensing range $R_{\max}^{\mathrm{sens}}$ is the distance at
which $\gamma_{\mathrm{sens}}=\gamma_{\mathrm{th}}$.

\subsection{Communication Branch}
\label{sec:comm_branch}

The one-way received signal at the remote receiver is
$r_c(t)=\sqrt{P_{\mathrm{rx}}(R)}\,x(t-\tau_c)$, with the received
power given by the free-space propagation law \cite{skolnik_introduction_1980},
\begin{equation}
P_{\mathrm{rx}}(R)=\frac{P_tG_tG_c\lambda^2}{(4\pi)^2R^2},
\label{eq:prx}
\end{equation}
and the same square-law detection applies. The received carrier now
plays the part the leakage plays in the sensing branch: it is the
reference the payload beats against, and it beats against the receiver
noise in the same way. There is no cross-beat, since there is only one
path, so
\begin{equation}
\gamma_{\mathrm{comm}}
=\frac{2m^2(1-\rho)P_{\mathrm{rx}}^2}
{N_{\mathrm{d},0}+2P_{\mathrm{rx}}N+\kappa m^4P_{\mathrm{rx}}^2},
\label{eq:sinr_comm}
\end{equation}
which is~\eqref{eq:sinr_sens} with the leakage replaced by the received
carrier. The desired term scales as $P_{\mathrm{rx}}^2\propto R^{-4}$,
matching the SI-homodyned sensing branch. At the ranges of interest the
carrier--noise beat dominates the denominator, while the SSBI imposes a
ceiling $2(1-\rho)/\kappa m^2$ that lies well above $\gamma_c$ for the
values of Table~\ref{tab:simulation_parameters}. The maximum
communication range $R_{\max}^{\mathrm{comm}}$ is the distance at which
$\gamma_{\mathrm{comm}}=\gamma_c$.

\subsection{ISAC Range Analysis}
\label{sec:isac_range}

A radar link and a communication link of the same platform reach
different distances, since one is governed by the two-way radar
equation and the other by the one-way propagation law, and their
respective maximum ranges must therefore be evaluated separately
\cite{sturm_waveform_2011}. A joint system delivers both functions
only where both are satisfied, so we define the ISAC range as
\begin{equation}
R_{\max}=\min\big(R_{\max}^{\mathrm{comm}},\,
R_{\max}^{\mathrm{sens}}\big).
\label{eq:rmax}
\end{equation}
The two thresholds behind it are of different nature. The
communication threshold $\gamma_c$ is a demodulation requirement: the
SINR at which the recovered constellation reaches the target error
rate, here $32$-QAM at the $20\%$-overhead SD-FEC threshold of
$\mathrm{BER}=2\times10^{-2}$, giving $\gamma_c=15.75$~dB. The sensing
threshold $\gamma_{\mathrm{th}}$ is a detection requirement: the SINR
at which the range-profile peak, after the coherent processing gain
$G_p$, is declared a target at a prescribed detection and false-alarm
probability, here $P_{\mathrm{d}}=0.9$ at $P_{\mathrm{fa}}=10^{-6}$,
giving $\gamma_{\mathrm{th}}=13.20$~dB from the Marcum
$Q$-function. Both are evaluated with the parameters of
Table~\ref{tab:simulation_parameters}.

Fig.~\ref{fig:isac_range} evaluates~\eqref{eq:rmax} against
$\sigma_{\mathrm{eff}}$ at $B=20$~GHz, whose $7.5$-mm range resolution
follows from $c/2B$. The communication range does not depend on the
target and is flat, whereas the sensing range grows with
$\sigma_{\mathrm{eff}}$. Both sensing curves scale as
$\sigma_{\mathrm{eff}}^{1/4}$, since $\gamma_{\mathrm{sens}}$ is
proportional to $\sigma_{\mathrm{eff}}R^{-4}$ with the leakage and to
$\sigma_{\mathrm{eff}}^2R^{-8}$ without it, so admitting the leakage
shifts the sensing range up by a constant factor.
Fig.~\ref{fig:isac_range}(a), with $\rho=0.2$, is
communication-centric: the maximum communication range is $1.1$~m, and
for a strong target the sensing range already exceeds it.
Fig.~\ref{fig:isac_range}(b), with $\rho=0.8$, is sensing-centric: the
maximum communication range falls to $0.8$~m, but a weak,
low-$\sigma_{\mathrm{eff}}$ target remains detectable.

\begin{figure}[!t]
\centering
\subfloat[$\rho=0.2$]{\includegraphics[width=0.48\columnwidth]{fig/fig3(a).png}}
\hfil
\subfloat[$\rho=0.8$]{\includegraphics[width=0.48\columnwidth]{fig/fig3(b).png}}
\caption{ISAC range versus effective radar cross-section for
(a)~$\rho=0.2$ and (b)~$\rho=0.8$. Both sensing curves scale as
$\sigma_{\mathrm{eff}}^{1/4}$; admitting the leakage shifts the sensing
range up by a constant factor.}
\label{fig:isac_range}
\end{figure}

Fig.~\ref{fig:sinr_si} shows the sensing SINR
of~\eqref{eq:sinr_sens} as a function of the leakage power at
$R=1.1$~m, and separates into three regions.

Where the leakage does not exceed the echo, the echo self-beat
dominates and the SINR approaches the leakage-free floor
of~\eqref{eq:sinr_nosi}, which lies below $\gamma_{\mathrm{th}}$: with
no leakage to homodyne the echo, the sensing threshold is not met.

Once $P_{\mathrm{SI}}$ exceeds $P_{\mathrm{ec}}$ the cross-beat takes
over. With only the fixed floor in the denominator,
\begin{equation}
\gamma_{\mathrm{sens}}\simeq
\rho G_p\,\frac{2m^2P_{\mathrm{SI}}P_{\mathrm{ec}}}{N_{\mathrm{d},0}},
\label{eq:sinr_window}
\end{equation}
the dashed line in Fig.~\ref{fig:sinr_si}, along which the SINR would
rise one for one with the leakage and gain the full
$P_{\mathrm{SI}}/P_{\mathrm{ec}}=17$~dB at the operating point.

The rise is curtailed from two directions. The SI--noise beat overtakes
the floor first and, growing at the same rate as the ranging term,
flattens the curve; the SSBI follows and, growing twice as fast, bends
it back down. The realized gain at the operating point is $12.5$~dB
rather than the ideal $17$~dB. The shaded
band marks the leakage powers for which the sensing threshold is met
without compressing the front end, and the isolation of
Table~\ref{tab:simulation_parameters} places the operating point inside
it. What matters is not that a single optimum exists, but that the gain
is available across a broad and easily met range of isolation.

\begin{figure}[!t]
\centering
\includegraphics[width=\columnwidth]{fig/fig4.png}
\caption{Sensing SINR versus self-interference power at the
communication limit of $1.1$~m. The shaded band marks the leakage
powers for which the sensing threshold is met without compressing the
front end.}
\label{fig:sinr_si}
\end{figure}

In summary, subcarrier intensity modulation enables simultaneous
communication and sensing with a single free-running photonic
transmitter and two square-law receivers, without a local oscillator.
Reusing the duplexer leakage as both a self-homodyne and a zero-delay
reference relaxes the sensing SINR from $R^{-8}$ to $R^{-4}$ and yields
absolute, drift-immune ranging through normalized CFR estimation. The
duplexer isolation is therefore a parameter to place within the
operating window of Fig.~\ref{fig:sinr_si} rather than one to maximize.

\begin{table}[t]
\caption{System parameters for ISAC range analysis}
\label{tab:simulation_parameters}
\centering
\renewcommand{\arraystretch}{1.08}
\begin{tabular}{@{}llr@{}}
\toprule
Parameter & Symbol & Value \\
\midrule
THz carrier frequency & $f_c$ & $280$ GHz \\
Intermediate frequency & $f_{\mathrm{IF}}$ & $11$ GHz \\
Waveform / modulation & -- & DFT-s-OFDM / 32-QAM \\
Bandwidth (symbol rate) & $B$ & $20$ GHz \\
Data rate & $R_b$ & $100$ Gb/s \\
Range resolution & $\Delta R=c/2B$ & $7.5$ mm \\
Pilot length & $N_p$ & $1024$ symbols \\
Processing gain & $G_p=N_p$ & $30.10$ dB \\
Power-allocation ratio & $\rho$ & $0.20$ \\
Modulation index & $m$ & $0.22$ \\
Effective in-band SSBI fraction$^{\S}$ & $\kappa$ & $0.069$ \\
Laser linewidth & $\Delta\nu$ & $15$ kHz \\
Transmitted THz power & $P_t$ & $-10$ dBm \\
Tx/Rx antenna gains & $G_t,G_r$ & $33$ dBi \\
Duplexer isolation & $\mathrm{ISO}$ & $25$ dB \\
SI power at LNA input & $P_{\mathrm{SI}}$ & $-35$ dBm \\
Echo power at $R_{\max}$ & $P_{\mathrm{ec}}$ & $-52$ dBm \\
LNA compression point & $\mathrm{IP}_{1\mathrm{dB}}$ & $-20$ dBm \\
Effective RCS$^{\P}$ & $\sigma_{\mathrm{eff}}$ & $-8$ dBsm \\
System noise figure & $F$ & $8$ dB \\
Noise power & $N$ & $-63.0$ dBm \\
Detector-output noise floor & $N_{\mathrm{d},0}$ & $-97.5$ dB(mW$^2$) \\
Required comm.\ SINR$^{\dagger}$ & $\gamma_c$ & $15.75$ dB \\
Required sensing SINR$^{\ddagger}$ & $\gamma_{\mathrm{th}}$ & $13.20$ dB \\
ISAC range & $R_{\max}$ & $1.1$ m \\
\bottomrule
\multicolumn{3}{@{}l}{\footnotesize $^{\P}$Nominal value; swept in Fig.~\ref{fig:isac_range}.} \\
\multicolumn{3}{@{}l}{\footnotesize $^{\S}$Conservative upper bound, obtained by assigning the whole measured non-AWGN excess to SSBI.} \\
\multicolumn{3}{@{}l}{\footnotesize $^{\dagger}$32-QAM at the $20\%$-overhead SD-FEC threshold, $\mathrm{BER}=2\times10^{-2}$.} \\
\multicolumn{3}{@{}l}{\footnotesize $^{\ddagger}P_{\mathrm{d}}=0.9$ at $P_{\mathrm{fa}}=10^{-6}$.} \\
\end{tabular}
\end{table}

\bibliographystyle{IEEEtran}
\bibliography{references}

\end{document}
